% WS2: literature survey organised around the final design (replaces the mid-term six-paradigm list).
% References checked against Crossref metadata (see fixes/WS2/reference_audit.md).
\section{Literature Survey and Positioning of the Final Design}
\label{sec:literature_survey}
\begin{finalpart}
This survey is organised by the building blocks of the final design (Sec.~\ref{sec:final_algo}) and states for each whether the design uses, rejects or only compares with it. Published accuracies on the same data are compared in Sec.~\ref{sec:v7_literature}.

\begin{enumerate}
    \item \textbf{Maternal template subtraction (used).} Subtracting an averaged maternal beat from each lead is the classic non-invasive method; Behar \etal~\cite{behar2014b} benchmark template-subtraction variants (TS, TSpca) and source-separation methods on CinC 2013 set-a, and Andreotti \etal~\cite{andreotti2014} adapt the template to each beat. The final design subtracts, per lead and per maternal beat, a fit on the mean and the top two principal components of the last 20 beats (Sec.~\ref{sec:final_algo}, step 3); its expected residual is analysed in Sec.~\ref{sec:analysis}.
    \item \textbf{Periodic component analysis and GEVD (used).} Sameni \etal~\cite{sameni2008} find spatial filters that maximise the periodic (beat-locked) structure of multichannel ECG with a generalised eigenvalue decomposition. The final design uses the same principle on its own detected beats, without labels, to compute 8 spatial filters that feed 12 of its 16 candidate tracks (Sec.~\ref{sec:final_algo}, step 4; derivation in Sec.~\ref{sec:analysis}).
    \item \textbf{Blind source separation, ICA (compared, not used).} ICA methods (JADE~\cite{cardoso1999}, FastICA~\cite{hyvarinen2000}, subspace separation~\cite{delathauwer2000}) give the highest published F1 on set-a when combined with template subtraction (TSpca-ICA and fusion, Tab.~\ref{tab:v7_literature}), and block-on-line JADE has run in real time on a floating-point DSP~\cite{muceli2007}. We did not use ICA and made no measurement comparing it with our candidate tracks, so we make no claim about its suitability for this FPGA; real-time ICA itself has been shown on embedded hardware~\cite{muceli2007}.
    \item \textbf{CinC 2013 challenge methods (compared).} The PhysioNet/CinC 2013 challenge~\cite{silva2013} defined set-a and the $\pm 50\,$ms beat scoring that we use. Its entries and later work~\cite{behar2014b,andreotti2014} run offline on whole records; Sec.~\ref{sec:v7_literature} compares them with v7-DL under that caveat.
    \item \textbf{Beat tracking and track selection (our design).} The final design keeps 16 candidate tracks with RR-predictive trackers and chooses among them with a label-free quality score, causally (v6, v6.1) or with a fixed look-ahead and a Viterbi step (v7-DL, Sec.~\ref{sec:v7dl}). We did not find a published fetal ECG method with this exact selection scheme and do not claim to compare with one.
    \item \textbf{Adaptive filtering with a maternal reference (rejected).} Adaptive noise cancelling~\cite{widrow1975} needs a maternal reference lead; Edwin Dhas and Suchetha~\cite{edwin2022} use abdominal and thoracic ECG with a non-causal adaptive filter. The final design has no thoracic lead (CinC set-a and ADFECGDB provide only abdominal leads) and is causal.
    \item \textbf{Embedded and FPGA fetal ECG (compared).} Vasudeva \etal~\cite{vasudeva2020} implement preprocessing, LMS adaptive-filter extraction and FHR detection on an Artix-7 FPGA; Muceli \etal~\cite{muceli2007} run JADE on a DSP. Our design differs in running the whole multi-lead algorithm with 16 candidate tracks as RTL on a Spartan-7 with a PC making five low-rate decisions (Sec.~\ref{sec:final_hw}); it uses 43 DSP slices, not a single DSP.
    \item \textbf{Mid-term paradigms (superseded).} Hankel SVD~\cite{kanjilal1997,callaerts1990}, nonlinear Bayesian filtering~\cite{sameni2007} and the 1D U-Net~\cite{ronneberger2015} were part of the mid-term comparison (App.~\ref{app:midterm}); the U-Net's real-data result is in Sec.~\ref{sec:final_realdata}.
\end{enumerate}
\end{finalpart}
